\section{Worked examples}
\label{sec:worked-examples}

An edit makes a rule easier to use by exposing the change. This section follows
a real revision from draft to teaching rewrite and then explains the editorial
decisions. Shreyas prepared the source notes for the lab. Because the notes
contain feedback rather than a finished replacement, the ``after'' passage is
a reconstruction. Its paper-specific claims still need to be checked against
the manuscript and experiments.

\subsection{Introduction}

\begin{AIbox}[breakable]{Before}
A long-standing goal in robotics is to build agents that can improve
autonomously: given a new task, a robot should attempt it, learn from its own
experience, and rapidly become competent without requiring large amounts of
additional expert data. A natural milestone toward this goal is few-shot
adaptation, where a small number of demonstrations specifies the task and the
robot uses subsequent interaction to refine its behavior. This paradigm has
been highly effective in language, where large pretrained models can adapt to
new tasks from only a few examples in context.

Prior work has approached this problem from several directions, but none
directly resolves the minimal-data adaptation setting. One line of work studies
in-context or prompt-conditioned robot policies, where demonstrations, goals,
or language instructions are provided as part of the input to specify a new
task. While these methods make task specification more flexible, they typically
rely on large-scale imitation training over many tasks and do not use robot
interaction to autonomously improve after receiving a new demonstration. A
second line of work pretrains generalist robot policies on large multi-task
datasets and then finetunes them to new domains, showing that broad pretraining
can provide useful visual and motor priors. However, these approaches still
leave open how to adapt reliably when only one or a few downstream
demonstrations are available, and how to autonomously improve via
robot collected data. Finally, offline-RL-based adaptation methods such as PTR
show that task-specific demonstrations can be combined with broad prior robot
data to learn new tasks, but still require multiple demonstrations and
primarily treat demonstrations as data for learning the downstream policy.
\end{AIbox}

\begin{AIbox}[breakable]{After: teaching reconstruction}
A long-standing goal in robotics is to build robots that improve autonomously.
Given a new task, a robot should learn from its own interaction and become
competent with little additional expert data. Pretrained robot policies provide
stronger starting points, but adapting them to a new task often assumes access
to many expert demonstrations. We study a stricter regime in which expert data
is scarce but autonomous interaction remains available: given a pretrained
policy and one to five demonstrations, when can a robot improve through its own
experience?

The central hypothesis is that the demonstrations need not cover every
successful behavior. Instead, they anchor the pretrained policy near
task-relevant behavior. Adapting the action head provides a warm start, and
online reinforcement learning can then learn corrections through interaction.
The hypothesis separates three empirical questions: the information retained
by the pretrained representation, the behavior covered by demonstrations, and
the corrections learned after the warm start.

This setting complements in-context, imitation-based, and offline adaptation
rather than competing with them. Those approaches use demonstrations or prior
data to specify behavior; here, autonomous interaction remains available after
demonstrations become the bottleneck. The paper therefore defines a constrained
adaptation problem and tests a minimal recipe for it. The constraint also
changes the evaluation. Initial success may not predict adaptability, so the
evaluation should measure improvement from limited human input.
\end{AIbox}

\subsubsection*{The revision starts with the bottleneck}

The rewrite opens on autonomous improvement and then names scarce
demonstrations as the bottleneck. In-context learning no longer carries the
motivation. The research question fixes the base policy, demonstration budget,
interaction resource, and desired improvement before the literature appears.

Prior work now enters through those choices. In-context and imitation-based
methods specify behavior, while online reinforcement learning uses subsequent
interaction. The paper can therefore complement those methods without claiming
that they fail at their own objective.

The middle paragraph explains the effect of each ingredient. Pretraining
supplies the representation, demonstrations anchor task-relevant behavior,
action adaptation supplies a warm start, and online learning tests local
correction. Experiments can challenge each part of that account.

The rewrite also matches the claim to the study. The paper defines a
constrained problem and tests a recipe; it does not solve adaptation. This
version removes the ``prior work / prior data'' echo, shortens the catalogue of
approaches, and lets each paragraph establish one inference for the reader.

\subsubsection*{The frame leads to four experimental comparisons}

Base-policy experiments should measure the content and alignment of the
pretrained representation. A separate comparison should isolate the benefit of
updating the visual-language backbone against action-head-only adaptation.

Demonstration experiments should vary example count and task-mode coverage
separately. The notes propose that one to three examples per covered mode may
suffice and that broader mode coverage may matter more than repeated examples
of one mode. These remain hypotheses until the experiments establish them.

Algorithm experiments should isolate the gain from online reinforcement
learning after the warm start. Ablations can test aggressive critic updates
and a success buffer, then measure residual correction beyond behavior
cloning. Offline-to-online alternatives, inference-time steering or reasoning,
and WAM hardware results belong only if the paper includes those tests.

The evaluation should distinguish in-distribution from out-of-distribution
adaptation, position from visual shifts, and initial success from eventual
fine-tunability. The source notes propose position shifts of 5, 10, and
20\,cm in LIBERO and Robocasa alongside a standard visual-perturbation setting.
These details belong in the paper only after the authors verify the protocol.

The contribution order should follow the evidence. Define the constrained
problem and test the minimal recipe first. Then report the necessary conditions
and the limits of imitation and offline-to-online methods in this regime. Close
with the consequences for evaluation. Each contribution should state the
outcome of the study. A planned ablation or broad hope is not yet a finding.

The source notes end with a one-sentence working hypothesis:

\begin{quote}
Minimal demonstrations do not make adaptation possible by covering the task;
they make adaptation possible by anchoring a pretrained policy in the right
region, where autonomous online reinforcement learning can finish the work.
\end{quote}

That sentence gives the paper a mechanism to test. It should remain a
hypothesis until the representation, coverage, and adaptation experiments
support each part.
